\documentclass[11pt,a4paper]{article}

% ── Packages ──────────────────────────────────────────────
\usepackage[utf8]{inputenc}
\usepackage[T1]{fontenc}
\usepackage[margin=1in]{geometry}
\usepackage{hyperref}
\usepackage{xcolor}
\usepackage{listings}
\usepackage{booktabs}
\usepackage{longtable}
\usepackage{array}
\usepackage{enumitem}
\usepackage{fancyhdr}
\usepackage{titlesec}
\usepackage{caption}
\usepackage{mdframed}

% ── Colors & Hyperref ─────────────────────────────────────
\definecolor{codebackground}{HTML}{F5F5F5}
\definecolor{linkcolor}{HTML}{2563EB}
\definecolor{notebackground}{HTML}{EEF2FF}
\definecolor{noteborder}{HTML}{6366F1}
\definecolor{placeholder}{HTML}{DC2626}

\hypersetup{
    colorlinks=true,
    linkcolor=linkcolor,
    urlcolor=linkcolor,
    pdftitle={MrBERT -- First Milestone Run Plan},
}

% ── Code listing style ────────────────────────────────────
\lstdefinestyle{bashstyle}{
    backgroundcolor=\color{codebackground},
    basicstyle=\ttfamily\small,
    breaklines=true,
    keepspaces=true,
    frame=single,
    framerule=0.4pt,
    rulecolor=\color{black!20},
    xleftmargin=1em,
    xrightmargin=0.5em,
    aboveskip=1em,
    belowskip=0.8em,
    language=bash,
}

\lstset{style=bashstyle}

% ── Note/tip box ──────────────────────────────────────────
\mdfdefinestyle{notestyle}{
    backgroundcolor=notebackground,
    linecolor=noteborder,
    linewidth=1.5pt,
    leftline=true,
    rightline=false,
    topline=false,
    bottomline=false,
    innerleftmargin=10pt,
    innerrightmargin=10pt,
    innertopmargin=8pt,
    innerbottommargin=8pt,
    skipabove=0.8em,
    skipbelow=0.8em,
}

% ── Header/Footer ─────────────────────────────────────────
\pagestyle{fancy}
\fancyhf{}
\fancyhead[L]{\small\leftmark}
\fancyhead[R]{\small MrBERT Milestone Run Plan}
\fancyfoot[C]{\thepage}
\renewcommand{\headrulewidth}{0.4pt}

% ── Column types ──────────────────────────────────────────
\newcolumntype{L}[1]{>{\raggedright\arraybackslash}p{#1}}
\newcolumntype{C}[1]{>{\centering\arraybackslash}p{#1}}

% ── Placeholder command ───────────────────────────────────
\newcommand{\tbd}{\textcolor{placeholder}{\texttt{\_\_\_}}}

% ══════════════════════════════════════════════════════════
\begin{document}

% ── Title ─────────────────────────────────────────────────
\begin{center}
    {\LARGE \textbf{MrBERT --- First Milestone Run Plan}}\\[1em]
\end{center}

\textbf{Goal:} Produce a complete set of training runs, metrics, and charts that tell a coherent story at the advisor meeting: \emph{``MrBERT learns to delete tokens meaningfully. At 30\% deletion, accuracy drops minimally while inference gets measurably faster. Random deletion at the same rate performs worse, confirming the gate is learning something non-trivial.''}

\tableofcontents
\bigskip


% ══════════════════════════════════════════════════════════
\section{Run Index}
\label{sec:run-index}

\begin{table}[h]
\centering
\small
\begin{tabular}{@{}cllclllc@{}}
\toprule
\textbf{ID} & \textbf{Model} & \textbf{Task} & \textbf{Del} & \textbf{Variant} & \textbf{W\&B Project} & \textbf{Run Name} & \textbf{Tier} \\
\midrule
A & BERT   & SNLI  & ---  & baseline                      & mrbert-snli  & bert-snli-baseline    & 1 \\
B & MrBERT & SNLI  & 0\%  & no deletion pressure          & mrbert-snli  & mrbert-snli-0pct      & 1 \\
C & MrBERT & SNLI  & 30\% & main result (soft del only)   & mrbert-snli  & mrbert-snli-30pct     & 1 \\
M & MrBERT & SNLI  & 30\% & hard\_delete\_train\_prob=0.5 & mrbert-snli  & mrbert-snli-30pct-hd  & 1 \\
D & MrBERT & SNLI  & 50\% & tradeoff curve                & mrbert-snli  & mrbert-snli-50pct     & 1 \\
E & MrBERT & SNLI  & 70\% & tradeoff curve                & mrbert-snli  & mrbert-snli-70pct     & 1 \\
F & MrBERT & SNLI  & 30\% & random gate                   & mrbert-snli  & mrbert-snli-random30  & 1 \\
G & MrBERT & SNLI  & 30\% & no PI controller              & mrbert-snli  & mrbert-snli-nopi      & 2 \\
H & MrBERT & SNLI  & 30\% & gate at layer 1               & mrbert-snli  & mrbert-snli-layer1    & 2 \\
I & MrBERT & SNLI  & 30\% & gate at layer 6               & mrbert-snli  & mrbert-snli-layer6    & 2 \\
J & MrBERT & SNLI  & 30\% & gate at layer 9               & mrbert-snli  & mrbert-snli-layer9    & 2 \\
K & BERT   & SQuAD & ---  & baseline                      & mrbert-squad & bert-squad-baseline   & 1 \\
L & MrBERT & SQuAD & 30\% & main result                   & mrbert-squad & mrbert-squad-30pct    & 1 \\
\bottomrule
\end{tabular}
\caption{Complete run index. \textbf{Tier~1} = essential for the advisor meeting. \textbf{Tier~2} = adds depth for the gate layer ablation and PI controller ablation; launch alongside Tier~1 if bandwidth allows.}
\label{tab:run-index}
\end{table}


% ══════════════════════════════════════════════════════════
\section{Phase 3 --- Capture Metrics from W\&B}
\label{sec:phase3}

\begin{mdframed}[style=notestyle]
\textbf{Sequential.} Do this while waiting for downloads or immediately after. These numbers will be substituted into the chart commands in Phase~6.
\end{mdframed}

For each completed run, open W\&B and copy the following values from the \textbf{Summary} tab. Record them in the tables below.

\subsection{SNLI Results Table}

\begin{table}[h]
\centering
\small
\begin{tabular}{@{}clccccc@{}}
\toprule
\textbf{Run} & \textbf{Model} & \textbf{Del Target} & \textbf{test/acc} & \textbf{\% non-pad del} & \textbf{new\_seq\_len} & \textbf{seq $\Delta$\%} \\
\midrule
A & BERT baseline             & 0\%  & 0.9048 & 0       & 128.0  & 0     \\
B & MrBERT 0\%                & 0\%  & 0.9050 & 0       & 27.41  & 94.65 \\
C & MrBERT 30\%               & 30\% & 0.9021 & 30.80   & 18.96  & 96.30 \\
M & MrBERT 30\% hard-train    & 30\% & 0.9026 & 31.14   & 18.87  & 96.31 \\
D & MrBERT 50\%               & 50\% & 0.8952 & 52.55   & 12.97  & 97.47 \\
E & MrBERT 70\%               & 70\% & 0.8825 & 72.43   & 7.53   & 98.53 \\
F & MrBERT Random 30\%        & 30\% & 0.8726 & 49.64   & 64.21  & 87.46 \\
G & MrBERT No-PI 30\%         & 30\% & 0.8647 & 88.98   & 3.00   & 99.41 \\
H & MrBERT Layer 1            & 30\% & 0.9042 & 30.98   & 18.92  & 96.30 \\
C & MrBERT Layer 3 (reuse C)  & 30\% & \multicolumn{4}{c}{\emph{(same as Run C above)}} \\
I & MrBERT Layer 6            & 30\% & 0.9036 & 33.89   & 18.07  & 96.47 \\
J & MrBERT Layer 9            & 30\% & \tbd   & \tbd    & \tbd   & \tbd  \\
\bottomrule
\end{tabular}
\caption{SNLI results from W\&B. Run J (Layer 9) is pending.}
\label{tab:snli-results}
\end{table}

\subsection{SQuAD Results Table}

\begin{table}[h]
\centering
\begin{tabular}{@{}clccccc@{}}
\toprule
\textbf{Run} & \textbf{Model} & \textbf{Del Rate} & \textbf{test/squad\_em} & \textbf{test/squad\_f1} & \textbf{seq $\Delta$\%} \\
\midrule
K & BERT baseline & 0\%  & \tbd & \tbd & 0    \\
L & MrBERT 30\%   & 30\% & \tbd & \tbd & \tbd \\
\bottomrule
\end{tabular}
\caption{SQuAD results (pending).}
\label{tab:squad-results}
\end{table}

\begin{mdframed}[style=notestyle]
\textbf{W\&B tip:} In each run's Summary tab, search for \texttt{test/} to find all final test metrics. The metrics \texttt{seq\_len\_reduction\_pct} and \texttt{new\_seq\_len} are under the eval prefix in the last logged step.
\end{mdframed}


% ══════════════════════════════════════════════════════════
\section{Phase 4 --- Runtime Measurement}
\label{sec:phase4}

\begin{mdframed}[style=notestyle]
\textbf{Run on A100 via Modal} --- do not run locally (CPU timing is unreliable for speedup claims). Requires all Phase~1 training runs to have completed (checkpoints in the Modal volume). Hard deletion is used by default: MrBERT tokens are physically removed from tensors on GPU.
\end{mdframed}

Launch the benchmark on A100 (uses all checkpoints already in the Modal volume):

\begin{lstlisting}
modal run training/train_modal.py::benchmark_main
\end{lstlisting}

This benchmarks all SNLI runs and gate-layer ablation runs against each other, then saves results to the volume. Download the results:

\begin{lstlisting}
modal volume get mrbert-checkpoints analysis_figures ./analysis/figures
\end{lstlisting}

\textbf{Fill in the table from \texttt{analysis/figures/runtime\_table.csv}:}

\begin{table}[h]
\centering
\begin{tabular}{@{}lcc@{}}
\toprule
\textbf{Model} & \textbf{ms/sample} & \textbf{\% decrease vs BERT} \\
\midrule
BERT baseline    & \tbd & ---  \\
MrBERT 0\%       & \tbd & \tbd \\
MrBERT 30\%      & \tbd & \tbd \\
MrBERT 30\% hd   & \tbd & \tbd \\
MrBERT 50\%      & \tbd & \tbd \\
MrBERT 70\%      & \tbd & \tbd \\
Random 30\%      & \tbd & \tbd \\
Layer 1          & \tbd & \tbd \\
Layer 3          & \tbd & \tbd \\
Layer 6          & \tbd & \tbd \\
Layer 9          & \tbd & \tbd \\
\bottomrule
\end{tabular}
\caption{Runtime measurements (fill in from \texttt{runtime\_table.csv}).}
\label{tab:runtime}
\end{table}

Saved files: \texttt{analysis/figures/runtime\_table.csv}, \texttt{analysis/figures/runtime\_vs\_deletion.pdf}


% ══════════════════════════════════════════════════════════
\section{Phase 5 --- Deletion Pattern Analysis}
\label{sec:phase5}

\begin{mdframed}[style=notestyle]
\textbf{Parallel with Phase~4.} Requires Phase~2 (MrBERT 30\% SNLI checkpoint downloaded).
\end{mdframed}

\begin{lstlisting}
cd mrbert/
modal volume get mrbert-checkpoints \
  mrbert-snli-30pct/final ./local_checkpoints/mrbert-snli-30pct
\end{lstlisting}

\begin{lstlisting}
cd mrbert/
python analysis/get_deletion_patterns.py \
  --model_path ./local_checkpoints/mrbert-snli-30pct/final \
  --local_snli_dir ./snli_datasets \
  --sample_size 1000 \
  --output_dir ./analysis/deletion_patterns
\end{lstlisting}

\begin{lstlisting}
cd mrbert/
python analysis/deletion_pattern_analysis.py \
  --input_file ./analysis/deletion_patterns/mrbert-snli-30pct_test.json \
  --output_dir ./analysis/figures
\end{lstlisting}

Saved files: \texttt{mrbert-snli-30pct\_test\_by\_type.pdf}, \texttt{mrbert-snli-30pct\_test\_premise\_vs\_hyp.pdf}


% ══════════════════════════════════════════════════════════
\section{Phase 6 --- Generate Charts}
\label{sec:phase6}

\begin{mdframed}[style=notestyle]
\textbf{Sequential.} Requires Phase~3 (W\&B metrics captured) and Phase~4 (runtime measured). Replace all \tbd{} placeholders with actual values before running. Run from the \texttt{mrbert/} directory.
\end{mdframed}

\subsection{Theoretical Compute Savings + Accuracy vs Compute + Accuracy vs Seq-Length Reduction}

Replace each \tbd{} with the actual \texttt{test/accuracy} value from the W\&B table:

\begin{lstlisting}
python analysis/compute_savings.py \
  --output_dir analysis/figures \
  --runs "BERT,___,0.0" "MrBERT-0%,___,0.0" \
         "MrBERT-30%,___,0.30" "MrBERT-50%,___,0.50" \
         "MrBERT-70%,___,0.70" "Random-30%,___,0.30"
\end{lstlisting}

Saved files: \texttt{macs\_relative.pdf}, \texttt{macs\_by\_gate\_layer.pdf}, \texttt{accuracy\_vs\_compute.pdf}, \texttt{accuracy\_vs\_seq\_reduction.pdf} (all in \texttt{analysis/figures/}).

\subsection{Gate Layer Ablation Chart}

Replace each \tbd{} with accuracy from W\&B and ms/sample from Phase~4:

\begin{lstlisting}
python analysis/compute_savings.py \
  --output_dir analysis/figures \
  --gate-layer-runs "Layer 1,___,___,1" "Layer 3,___,___,3" \
                    "Layer 6,___,___,6" "Layer 9,___,___,9"
\end{lstlisting}

Saved file: \texttt{analysis/figures/gate\_layer\_ablation.pdf}


% ══════════════════════════════════════════════════════════
\section{Complete List of Output Files}
\label{sec:output-files}

\begin{table}[h]
\centering
\small
\begin{tabular}{@{}L{6.5cm}L{5cm}L{3cm}@{}}
\toprule
\textbf{File} & \textbf{What It Shows} & \textbf{Advisor Relevance} \\
\midrule
\texttt{accuracy\_vs\_seq\_reduction.pdf}             & Accuracy vs sequence length reduction \% --- core tradeoff              & Primary result \\
\texttt{accuracy\_vs\_compute.pdf}                    & Accuracy vs relative MACs --- efficiency frontier                       & Primary result \\
\texttt{runtime\_vs\_deletion.pdf}                    & Measured ms/sample per model                                            & Core claim: actual speedup \\
\texttt{runtime\_table.csv}                           & Runtime table for paper (mirrors MrT5 Table~3)                          & For paper table \\
\texttt{gate\_layer\_ablation.pdf}                    & Accuracy + runtime vs gate layer --- dual axis                          & Ablation \\
\texttt{macs\_relative.pdf}                           & Theoretical compute savings curve                                       & Background \\
\texttt{macs\_by\_gate\_layer.pdf}                    & Theoretical savings by gate layer                                       & Background \\
\texttt{mrbert-snli-30pct\_test\_by\_type.pdf}        & Deletion rate by token type (word/subword/punct)                        & Qualitative \\
\texttt{mrbert-snli-30pct\_test\_premise\_vs\_hyp.pdf}& Premise vs hypothesis deletion rate                                     & Qualitative \\
\texttt{mrbert-snli-30pct\_test.json}                 & Per-token gate decisions, 1000 examples                                 & For colored examples in slides \\
\bottomrule
\end{tabular}
\caption{Complete list of output files (all in \texttt{analysis/figures/} unless otherwise noted).}
\label{tab:output-files}
\end{table}


% ══════════════════════════════════════════════════════════
\section{What to Present to Your Advisor}
\label{sec:presentation}

\subsection*{Structure (30--40 min)}

\subsubsection*{1. Motivation (2 min)}

\begin{itemize}[nosep]
    \item BERT processes all tokens with equal compute regardless of informativeness.
    \item MrT5 showed learned token deletion works for encoder-decoder (T5) models.
    \item Research question: can we apply the same mechanism to BERT for discriminative NLU tasks?
\end{itemize}

\subsubsection*{2. Architecture (5 min)}

\begin{itemize}[nosep]
    \item Delete gate inserted after encoder layer 3.
    \item Gate = \texttt{LayerNorm $\rightarrow$ Linear $\rightarrow$ scaled sigmoid}, 2,305 additional parameters out of 110M.
    \item Soft deletion during training: large negative bias to attention scores of deleted tokens.
    \item PI controller dynamically adjusts deletion pressure $\alpha$ to hit target deletion rate $\delta$.
    \item Show: gate equation + diagram of where it sits in the encoder stack.
\end{itemize}

\subsubsection*{3. Core Results --- SNLI (5 min)}

Show the filled-in results table + \texttt{accuracy\_vs\_seq\_reduction.pdf}:

\begin{table}[h]
\centering
\begin{tabular}{@{}lcccc@{}}
\toprule
\textbf{Model} & \textbf{Accuracy} & \textbf{Del Rate} & \textbf{Seq $\Delta$} & \textbf{Runtime} \\
\midrule
BERT baseline       & \tbd & 0\%  & 0\%  & \tbd{} ms \\
MrBERT 0\% (sanity) & \tbd & 0\%  & 0\%  & \tbd{} ms \\
MrBERT 30\%         & \tbd & 30\% & \tbd & \tbd{} ms \\
MrBERT 50\%         & \tbd & 50\% & \tbd & \tbd{} ms \\
MrBERT 70\%         & \tbd & 70\% & \tbd & \tbd{} ms \\
Random gate 30\%    & \tbd & 30\% & 30\% & \tbd{} ms \\
\bottomrule
\end{tabular}
\caption{Core SNLI results for advisor presentation.}
\label{tab:advisor-snli}
\end{table}

Key talking points:
\begin{itemize}[nosep]
    \item MrBERT 0\% should match BERT --- validates the implementation is fair.
    \item Compare MrBERT 30\% accuracy vs BERT baseline --- how small is the gap?
    \item Compare MrBERT 30\% vs Random 30\% --- the learned gate should do better, proving the gate is learning.
    \item Show \texttt{runtime\_vs\_deletion.pdf} --- the actual measured speedup.
\end{itemize}

\subsubsection*{4. Core Results --- SQuAD (5 min)}

\begin{table}[h]
\centering
\begin{tabular}{@{}lccccc@{}}
\toprule
\textbf{Model} & \textbf{EM} & \textbf{F1} & \textbf{Seq $\Delta$} & \textbf{Runtime} \\
\midrule
BERT baseline & \tbd & \tbd & 0\%  & \tbd{} ms \\
MrBERT 30\%   & \tbd & \tbd & \tbd & \tbd{} ms \\
\bottomrule
\end{tabular}
\caption{Core SQuAD results for advisor presentation.}
\label{tab:advisor-squad}
\end{table}

Key talking point: QA requires extracting the exact answer span. The gate must preserve semantically critical context tokens. Does the accuracy hold up?

\subsubsection*{5. Ablations (5 min)}

\begin{itemize}[nosep]
    \item \texttt{gate\_layer\_ablation.pdf}: earlier gate $\rightarrow$ faster but less accurate; layer 3 is the sweet spot.
    \item No PI controller: what happens to actual deletion rate without adaptive control?
    \item Show \texttt{mrbert-snli-30pct\_test\_by\_type.pdf}: which token types get deleted? (function words vs content words)
    \item Show \texttt{mrbert-snli-30pct\_test\_premise\_vs\_hyp.pdf}: does the gate treat premise and hypothesis differently?
\end{itemize}

\subsubsection*{6. Discussion (10 min)}

\begin{itemize}[nosep]
    \item The efficiency frontier: how far down the tradeoff curve is acceptable for publication?
    \item Does the gate learn linguistically meaningful patterns?
    \item What should the next set of experiments be?
\end{itemize}


% ══════════════════════════════════════════════════════════
\section{Questions for Your Advisor}
\label{sec:questions}

These questions will directly shape the second milestone experiments.

\subsection{Scope and Tasks}

\begin{enumerate}
    \item Is SNLI + SQuAD sufficient as the evaluation suite, or do you recommend adding a third task (e.g., MNLI for generalization, CoNLL NER for token classification, GLUE for a broader benchmark)?
    \item MrT5 uses 15 languages to demonstrate language-specific compression rates. For a BERT-based paper, should we include multilingual BERT (mBERT) with multilingual NLI to replicate this finding?
\end{enumerate}

\subsection{Architecture and Design}

\begin{enumerate}[resume]
    \item We use soft deletion throughout training (attention masking) and could optionally use hard deletion (physically remove tokens) at test time. Is the hard vs soft deletion comparison important to include in the paper?
    \item The gate is currently at layer 3 (out of 12). MrT5 uses layer 3 out of 12 for the same reason. Should we justify this choice with a thorough ablation (layers 1--11), or is a sample of 4 layers (1, 3, 6, 9) sufficient?
    \item We initialize MrBERT from pretrained BERT weights and only add 2,305 gate parameters. Is this the right framing, or should we also experiment with training the gate from scratch to test whether pretraining helps?
\end{enumerate}

\subsection{Baselines}

\begin{enumerate}[resume]
    \item MrT5 compares against Boundary Predictor and Convolutional Pooling baselines. For a BERT paper, are there equivalent compression baselines we should include (e.g., SpAtten, TR-BERT, LTP)?
    \item How much emphasis should the random deletion baseline receive? Is it sufficient to show MrBERT $>$ random, or do reviewers expect a richer set of non-learned baselines?
\end{enumerate}

\subsection{Metrics and Claims}

\begin{enumerate}[resume]
    \item The runtime speedup we measure is for forward-pass inference on a single GPU. For a publication claim, should we also measure wall-clock speedup in a batched serving scenario (where hard deletion breaks batch uniformity)?
    \item MrT5 reports bits-per-byte (language modeling loss). Our equivalent is accuracy/EM/F1. Is there any value in also pre-training MrBERT on MLM before fine-tuning, to get a language-modeling-style efficiency comparison?
\end{enumerate}

\subsection{Paper Framing}

\begin{enumerate}[resume]
    \item How should we frame the contribution relative to MrT5? Options:
    \begin{enumerate}[label=(\alph*)]
        \item Direct adaptation of MrT5 to BERT.
        \item Independent extension showing the mechanism generalizes across architectures.
        \item Applied study showing practical efficiency gains for discriminative NLU.
    \end{enumerate}
    \item What is the target venue? (ACL, EMNLP, NAACL, workshop?) --- this affects how thorough the ablations need to be and how many tasks are required.
    \item Is there a specific result threshold that would make this ``publishable'' --- e.g., a minimum runtime speedup \% at a maximum accuracy drop \%?
\end{enumerate}


\end{document}